% Fuente única de clase: Recursos/2025-11-06.pdf, páginas 1--11.
En la clase 21 usamos el lema de separación para estudiar submódulos de módulos libres sobre un DIP. Obtenemos bases compatibles, resoluciones libres y la clasificación de los módulos finitamente generados mediante factores invariantes y divisores elementales. La unicidad se demuestra recuperando invariantes del propio módulo.

%%%%%%%%%%%%%%%%%%%%%%%%%%%%%%%%%%%%%%%%%%%%%%%%%%%%%%%%%%%%%%%%%%%%%%%%%%%%%%%
\section{Submódulos libres y bases compatibles}
%%%%%%%%%%%%%%%%%%%%%%%%%%%%%%%%%%%%%%%%%%%%%%%%%%%%%%%%%%%%%%%%%%%%%%%%%%%%%%%
%%%%%%%%%%%%%%%%%%%%%%%%%%%%%%%%%%%%%%%%%%%%%%%%%%%%%%%%%%%%%%%%%%%%%%%%%%%%%%%
\subsection{Libertad de los submódulos}
%%%%%%%%%%%%%%%%%%%%%%%%%%%%%%%%%%%%%%%%%%%%%%%%%%%%%%%%%%%%%%%%%%%%%%%%%%%%%%%
% P01--P02. P01 repite el lema VI.5.2, probado en la clase 20.
En toda esta clase $R$ es un dominio de ideales principales, salvo cuando se indique una hipótesis más general. La página 1 retoma el \cref{lema:C20-separacion}: si $0\ne M\subseteq F$, con $F$ libre finitamente generado, podemos encontrar $y=ax\ne0$ y escribir
\[
F=Rx\oplus F',\qquad M=Ry\oplus(M\cap F').
\]
Usaremos también la elección de $a$ mediante una imagen maximal del submódulo $M$ en $R$.

\begin{prop}\label{prop:C21-submodulo-libre}
Si $F$ es libre de rango finito sobre un DIP y $M\subseteq F$ es un submódulo, entonces $M$ es libre de rango a lo sumo $\rk_R F$.
\end{prop}
\begin{proof}
Sea $n=\rk_R F$. Si $M=0$, la base es vacía. En caso contrario, el lema da $y_1\ne0$ y $M^{(1)}\subseteq M$ tales que $M=Ry_1\oplus M^{(1)}$. Si $M^{(1)}\ne0$, aplicamos otra vez el lema a $M^{(1)}\subseteq F$, y obtenemos
\[
M=Ry_1\oplus Ry_2\oplus M^{(2)}.
\]
Continuamos mientras el último submódulo sea no nulo. Como $F$ es libre de torsión, cada $Ry_i$ es libre con base $y_i$; la suma directa asegura que los $y_i$ son independientes. La \cref{prop:C20-cardinal} impide obtener más de $n$ elementos independientes dentro de $F$. Si después de $n$ pasos quedara un submódulo no nulo, el lema produciría uno más, contradicción. Por tanto, en algún paso $m\leq n$ se tiene $M^{(m)}=0$, y
\[
M=Ry_1\oplus\cdots\oplus Ry_m.
\]
Los elementos $y_1,\ldots,y_m$ forman la base buscada.
\end{proof}
En esta aplicación se mantiene $F$ como ambiente en cada paso. Así se completa el argumento iterativo de la página 2 sin suponer de antemano que los complementos $F'$ ya son libres. El cero que allí se representa a veces por $\varnothing$ debe leerse como submódulo cero.

%%%%%%%%%%%%%%%%%%%%%%%%%%%%%%%%%%%%%%%%%%%%%%%%%%%%%%%%%%%%%%%%%%%%%%%%%%%%%%%
\subsection{Una base del ambiente adaptada al submódulo}
%%%%%%%%%%%%%%%%%%%%%%%%%%%%%%%%%%%%%%%%%%%%%%%%%%%%%%%%%%%%%%%%%%%%%%%%%%%%%%%
% P03--P04.
\begin{coro}\label{coro:C21-bases-compatibles}
Sea $M\subseteq F$, con $F$ libre de rango $n$ sobre un DIP. Existen una base $x_1,\ldots,x_n$ de $F$, un entero $0\leq m\leq n$ y elementos no nulos $a_1,\ldots,a_m$ tales que
\[
a_1x_1,\ldots,a_mx_m\text{ es una base de }M,
\qquad a_1\mid a_2\mid\cdots\mid a_m.
\]
\end{coro}
\begin{proof}
Procedemos por inducción sobre $n$, incluyendo $n=0$. Si $M=0$, cualquier base de $F$ sirve, con $m=0$. Para $M\ne0$, aplicamos el lema de separación y escribimos
\[
F=Rx_1\oplus F',\qquad M=Ra_1x_1\oplus M',\qquad
M'=M\cap F',
\]
donde $(a_1)=\alpha(M)$ fue elegido maximal entre las imágenes $\varphi(M)$ para $\varphi:F\to R$. Por la proposición, $F'\subseteq F$ es libre de rango finito. Si $E'$ es una base de $F'$, la familia $\{x_1\}\cup E'$ es una base de $F$, de modo que $|E'|=n-1$ por la invariancia del número de base. Esto justifica la disminución del rango que se usa en la página 3.

La inducción aplicada a $M'\subseteq F'$ da una base $x_2,\ldots,x_n$ y una base $a_2x_2,\ldots,a_mx_m$ de $M'$, con $a_2\mid\cdots\mid a_m$. Junto con $x_1$ y $a_1x_1$ obtenemos las bases requeridas. Falta, si $m\geq2$, probar $a_1\mid a_2$.

Definamos $\varphi:F\to R$ por $\varphi(x_1)=\varphi(x_2)=1$ y $\varphi(x_i)=0$ para $i>2$. Como $\varphi(a_1x_1)=a_1$,
\[
(a_1)\subseteq\varphi(M).
\]
La maximalidad da igualdad. Entonces $a_2=\varphi(a_2x_2)\in(a_1)$, es decir, $a_1\mid a_2$, como se quería.
\end{proof}

%%%%%%%%%%%%%%%%%%%%%%%%%%%%%%%%%%%%%%%%%%%%%%%%%%%%%%%%%%%%%%%%%%%%%%%%%%%%%%%
\subsection{Caracterización mediante resoluciones libres}
%%%%%%%%%%%%%%%%%%%%%%%%%%%%%%%%%%%%%%%%%%%%%%%%%%%%%%%%%%%%%%%%%%%%%%%%%%%%%%%
% P04.
\begin{prop}\label{prop:C21-resoluciones}
Un dominio $R$ es DIP si y sólo si, para todo módulo finitamente generado $M$ y toda sobreyección $\pi_0:R^{m_0}\to M$, existe una sucesión exacta
\[
\begin{tikzcd}[column sep=large]
0\arrow[r]&R^{m_1}\arrow[r,"\pi_1"]&R^{m_0}\arrow[r,"\pi_0"]&M\arrow[r]&0
\end{tikzcd}
\]
con $m_1$ finito.
\end{prop}
\begin{proof}
Si $R$ es DIP, el núcleo de $\pi_0$ es libre de rango finito por la \cref{prop:C21-submodulo-libre}. Elegimos un isomorfismo $R^{m_1}\to\ker\pi_0$ y lo componemos con la inclusión en $R^{m_0}$. Así se obtiene $\pi_1$, inyectivo y con imagen exactamente $\ker\pi_0$. La sobreyectividad de $\pi_0$ asegura la exactitud restante.

A la inversa, sea $I$ un ideal cualquiera de $R$. El módulo $R/I$ es cíclico, luego finitamente generado. Aplicamos la hipótesis a la proyección $R\to R/I$. Su núcleo $I$ es isomorfo a un módulo libre $R^{m_1}$. Una base de $I\subseteq R$ es independiente en $R$, así que tiene a lo sumo un elemento por la \cref{prop:C20-cardinal}, válida para cualquier dominio. Por tanto $I=0$ o $I$ está generado por un solo elemento. Todos los ideales son principales.
\end{proof}

%%%%%%%%%%%%%%%%%%%%%%%%%%%%%%%%%%%%%%%%%%%%%%%%%%%%%%%%%%%%%%%%%%%%%%%%%%%%%%%
\section{Clasificación: existencia de las descomposiciones}
%%%%%%%%%%%%%%%%%%%%%%%%%%%%%%%%%%%%%%%%%%%%%%%%%%%%%%%%%%%%%%%%%%%%%%%%%%%%%%%
%%%%%%%%%%%%%%%%%%%%%%%%%%%%%%%%%%%%%%%%%%%%%%%%%%%%%%%%%%%%%%%%%%%%%%%%%%%%%%%
\subsection{Rango de un módulo finitamente generado}
%%%%%%%%%%%%%%%%%%%%%%%%%%%%%%%%%%%%%%%%%%%%%%%%%%%%%%%%%%%%%%%%%%%%%%%%%%%%%%%
% P05.
\begin{defi}
Para un módulo finitamente generado sobre un dominio, $\rk_R M$ es el máximo número de elementos linealmente independientes de $M$.
\end{defi}

\begin{ejer}[Remisión a VI.5.6]\label{ejer:C21-rango-finito}
Justificar que ese máximo existe y es finito. Comprobar que coincide con el rango ya definido cuando $M$ es libre de rango finito.
\end{ejer}
\begin{sol}
Sean $n$ generadores de $M$ y $\pi:R^n\to M$ la sobreyección inducida. Si $m_1,\ldots,m_t$ son independientes, elegimos levantamientos $v_i\in R^n$. Una relación $\sum r_iv_i=0$ produce $\sum r_im_i=0$, de modo que todos los $r_i$ son cero. Por tanto los $v_i$ son independientes y $t\leq n$ por la \cref{prop:C20-cardinal}. La misma cota impide una familia independiente infinita, pues contendría $n+1$ elementos. Los tamaños posibles pertenecen al conjunto finito $\{0,\ldots,n\}$, que incluye el tamaño cero, y por ello tienen un máximo. Si $M$ es libre, una base alcanza la cota y la misma proposición compara cualquier familia independiente con ella.
\end{sol}

%%%%%%%%%%%%%%%%%%%%%%%%%%%%%%%%%%%%%%%%%%%%%%%%%%%%%%%%%%%%%%%%%%%%%%%%%%%%%%%
\subsection{Factores invariantes y divisores elementales}
%%%%%%%%%%%%%%%%%%%%%%%%%%%%%%%%%%%%%%%%%%%%%%%%%%%%%%%%%%%%%%%%%%%%%%%%%%%%%%%
% P05--P07. La unicidad se demuestra en la sección siguiente.
\begin{ejer}[Comprobación remitida a la tarea 4]\label{ejer:C21-chino}
Para una familia finita de ideales comaximales dos a dos de un anillo conmutativo, demostrar que su producto es su intersección y que el cociente por ésta es el producto de sus cocientes.
\end{ejer}
\begin{sol}
Primero sean $I+J=R$. Escribamos $u+v=1$ con $u\in I$, $v\in J$. Siempre $IJ\subseteq I\cap J$. Si $z\in I\cap J$, entonces $z=zu+zv\in IJ$, luego hay igualdad. El mapa $R\to R/I\oplus R/J$ que toma las dos clases tiene núcleo $I\cap J$ y es sobreyectivo: para obtener $(a+I,b+J)$ usamos $av+bu$. El primer isomorfismo prueba el caso de dos ideales.

Para pasar a una familia finita, observemos que si cada $I_i$ es comaximal con $J$, podemos elegir $u_i\in I_i$ con $u_i\equiv1\pmod J$. Su producto pertenece a $\prod_iI_i$ y sigue siendo congruente con $1$ módulo $J$, por lo que $(\prod_iI_i)+J=R$. Aplicando sucesivamente el caso de dos ideales obtenemos tanto $\prod_iI_i=\bigcap_iI_i$ como el isomorfismo de cocientes. Para una familia de un elemento la afirmación es inmediata; para la vacía, el ideal es $R$ y el cociente y la suma de cocientes son cero.
\end{sol}
Se resuelven las afirmaciones efectivamente utilizadas por la remisión manuscrita de la página 7, sin suponer acceso al enunciado completo de esa tarea.

\begin{teo}[Clasificación sobre un DIP]\label{teo:C21-clasificacion}
Sea $M$ un módulo finitamente generado sobre un DIP $R$, y sea $r=\rk_R M$. Existen elementos no nulos y no unidades $a_1,\ldots,a_m$, con $a_1\mid\cdots\mid a_m$, tales que
\[
M\cong R^r\oplus R/(a_1)\oplus\cdots\oplus R/(a_m).
\]
Ésta es la descomposición en \emph{factores invariantes}. También existen primos no nulos $q_1,\ldots,q_t$, dos a dos no asociados, y exponentes positivos $e_{ij}$ tales que
\[
M\cong R^r\oplus\bigoplus_{i,j}R/(q_i^{e_{ij}}).
\]
Ésta es la descomposición en \emph{divisores elementales}. Las sumas son finitas y se permiten sumas vacías. Ambas descomposiciones son únicas: quedan determinados $r$, los ideales $(a_j)$ en su orden de divisibilidad y, para cada ideal primo $(q_i)$, el multiconjunto de exponentes.
\end{teo}
\begin{proof}[Existencia]
Elijamos $n$ generadores de $M$ y la sobreyección $\pi:R^n\to M$. Por el \cref{coro:C21-bases-compatibles}, existen una base $x_1,\ldots,x_n$ de $R^n$ y elementos no nulos $b_1\mid\cdots\mid b_s$ tales que $b_1x_1,\ldots,b_sx_s$ forman una base de $\ker\pi$. El primer teorema de isomorfismo da
\[
M\cong\frac{Rx_1\oplus\cdots\oplus Rx_n}
 {Rb_1x_1\oplus\cdots\oplus Rb_sx_s\oplus0\oplus\cdots\oplus0}.
\]
Para módulos $N_i\subseteq M_i$, el mapa coordenado
\[
\bigoplus_i M_i\longrightarrow\bigoplus_i M_i/N_i,
\qquad (m_i)_i\longmapsto(m_i+N_i)_i
\]
es sobreyectivo y tiene núcleo $\bigoplus_i N_i$. De aquí la identificación de cocientes de la página 6:
\[
M\cong R/(b_1)\oplus\cdots\oplus R/(b_s)\oplus R^{n-s}.
\]
Si algún $b_i$ es unidad, su cociente es cero y se omite. Las unidades forman un segmento inicial de la cadena, por lo que los factores restantes conservan la divisibilidad; los llamamos $a_1,\ldots,a_m$.

La suma $T=\bigoplus_jR/(a_j)$ es de torsión: el producto no nulo $h=\prod_j a_j$ la anula; si no hay factores, tomamos $h=1$. La parte $R^{n-s}$ es libre de torsión, así que $\operatorname{Tor}_R(M)$ corresponde precisamente a $0\oplus T$. Además, el rango de esta descomposición es $n-s$. En efecto, las coordenadas libres dan esa cantidad de vectores independientes. Si $m_1,\ldots,m_u$ son independientes en $M$, sus proyecciones a $R^{n-s}$ lo son: una relación entre esas proyecciones produciría $\sum c_im_i\in T$; al multiplicar por $h$ se obtiene $\sum hc_im_i=0$, luego $hc_i=0$ y $c_i=0$. Por tanto $u\leq n-s$ y concluimos $r=n-s$.

Para la segunda descomposición usamos que un DIP es DFU. Cada $a_j$ se factoriza, salvo unidad, como producto de potencias de primos distintos. Para un $j$ fijo, los ideales $(q_i^{e_{ij}})$ que aparecen son comaximales dos a dos, porque sus generadores tienen mcd igual a uno. Su producto coincide con su intersección y el teorema chino del resto da
\[
R/(a_j)\cong
R/\!\left(\bigcap_i(q_i^{e_{ij}})\right)
\cong\bigoplus_iR/(q_i^{e_{ij}}).
\]
Juntando los sumandos obtenemos los divisores elementales. Los primos que no dividen un $a_j$ simplemente no aportan un sumando.
\end{proof}
La justificación de que el rango es el tamaño de la parte libre se ha explicitado para completar la identificación usada en la página 7. La unicidad ocupa las siguientes subsecciones.

%%%%%%%%%%%%%%%%%%%%%%%%%%%%%%%%%%%%%%%%%%%%%%%%%%%%%%%%%%%%%%%%%%%%%%%%%%%%%%%
\section{Anuladores y unicidad}
%%%%%%%%%%%%%%%%%%%%%%%%%%%%%%%%%%%%%%%%%%%%%%%%%%%%%%%%%%%%%%%%%%%%%%%%%%%%%%%
%%%%%%%%%%%%%%%%%%%%%%%%%%%%%%%%%%%%%%%%%%%%%%%%%%%%%%%%%%%%%%%%%%%%%%%%%%%%%%%
\subsection{Módulos cíclicos y anuladores}
%%%%%%%%%%%%%%%%%%%%%%%%%%%%%%%%%%%%%%%%%%%%%%%%%%%%%%%%%%%%%%%%%%%%%%%%%%%%%%%
% P08--P09.
Las definiciones siguientes tienen sentido sobre cualquier anillo conmutativo.
\begin{defi}
Un módulo es \emph{cíclico} si está generado por un elemento. El \emph{anulador} de $M$ es
\[
\operatorname{Ann}_R(M)
=\set{a\in R\given am=0\text{ para todo }m\in M}.
\]
\end{defi}
Si $M=Rm$, el homomorfismo $R\to M$, $a\mapsto am$, es sobreyectivo; por el primer isomorfismo $M\cong R/I$, con $I$ su núcleo. Recíprocamente, $R/I$ está generado por $1+I$.

\begin{ejer}[Remisión a VI.4.4]\label{ejer:C21-anulador}
Probar que $\operatorname{Ann}_R(M)$ es un ideal. Si $R$ es un dominio y $M$ es finitamente generado, probar que $M$ es de torsión si y sólo si $\operatorname{Ann}_R(M)\ne(0)$.
\end{ejer}
\begin{sol}
El cero anula a $M$. Si $a,b$ lo anulan, entonces $(a-b)m=am-bm=0$ para todo $m$. Si $c\in R$, también $(ca)m=c(am)=0$. Esto verifica las condiciones de ideal.

Si $M$ está generado por $m_1,\ldots,m_u$ y es de torsión, elegimos $a_i\ne0$ con $a_im_i=0$. El producto $a=a_1\cdots a_u$ es no nulo porque $R$ es dominio y anula a todos los generadores, luego a todo $M$. Si la familia es vacía, $M=0$ y el producto vacío $1$ lo anula. En sentido contrario, un elemento no nulo del anulador sirve como testigo de torsión para cada elemento de $M$, incluso sin generación finita.
\end{sol}

\begin{lema}\label{lema:C21-anuladores}
Se tiene $\operatorname{Ann}_R(R/I)=I$ y
\[
\operatorname{Ann}_R(M_1\oplus M_2)
=\operatorname{Ann}_R(M_1)\cap\operatorname{Ann}_R(M_2).
\]
\end{lema}
\begin{proof}
Si $a$ anula a $R/I$, entonces $a(1+I)=0$, luego $a\in I$. Si $a\in I$, para todo $b$ se cumple $ab\in I$ y $a(b+I)=0$. Para la suma directa, un escalar anula todos los pares si y sólo si anula por separado los pares $(m_1,0)$ y $(0,m_2)$, que es exactamente la pertenencia a ambos anuladores. La misma prueba se aplica a sumas finitas por inducción.
\end{proof}

\begin{lema}\label{lema:C21-primos-anulador}
Sea $0\ne T$ un módulo de torsión finitamente generado sobre un DIP, expresado en las formas del \cref{teo:C21-clasificacion}. Entonces $\operatorname{Ann}_R(T)=(a_m)$, y los ideales primos que contienen ese anulador son exactamente los $(q_i)$ que aparecen en sus divisores elementales.
\end{lema}
\begin{proof}
Por el lema anterior y la cadena de divisibilidad,
\[
\operatorname{Ann}_R(T)=(a_1)\cap\cdots\cap(a_m)=(a_m).
\]
También se puede expresar esa intersección como el ideal del mcm, que aquí es $(a_m)$. Un primo $q_i$ que aparece en algún $a_j$ divide a $a_m$; por tanto $(a_m)\subseteq(q_i)$. Recíprocamente, si un ideal primo $\mathfrak p$ contiene a $(a_m)$, es no nulo. Al ser $R$ DIP, $\mathfrak p=(q)$ con $q$ primo no nulo. La inclusión equivale a $q\mid a_m$, así que $q$ es asociado a uno de los $q_i$ de la factorización.
\end{proof}
Para $T=0$, su anulador es $R$, ningún primo propio lo contiene y la lista elemental es vacía. La fórmula con $a_m$ se reserva para $T\ne0$, cuando existe un último factor.

%%%%%%%%%%%%%%%%%%%%%%%%%%%%%%%%%%%%%%%%%%%%%%%%%%%%%%%%%%%%%%%%%%%%%%%%%%%%%%%
\subsection{Recuperación de los exponentes}
%%%%%%%%%%%%%%%%%%%%%%%%%%%%%%%%%%%%%%%%%%%%%%%%%%%%%%%%%%%%%%%%%%%%%%%%%%%%%%%
% P07 y P10.
\begin{proof}[Unicidad en el \cref{teo:C21-clasificacion}]
El submódulo $T=\operatorname{Tor}_R(M)$ se define a partir de $M$ y todo isomorfismo lo preserva, junto con su anulador. El cociente $M/T\cong R^r$ recupera la parte libre; su rango y, por tanto, $r$ son intrínsecos. Por el \cref{lema:C21-primos-anulador}, también queda determinada la lista de ideales primos que intervienen en $T$.

Fijemos uno de sus generadores primos $q$. Consideremos el submódulo
\[
T_q=\set{m\in T\given q^um=0\text{ para algún entero }u\geq1}.
\]
Es cerrado bajo suma usando la mayor de dos potencias y bajo escalares por conmutatividad. Todo isomorfismo lleva $T_q$ sobre el submódulo correspondiente. En la descomposición elemental, $T_q$ es exactamente la suma de los factores de potencias de $q$: éstos quedan anulados por una potencia común. En un factor $R/(p^e)$ con $p$ no asociado a $q$, la identidad de Bézout para $q^u$ y $p^e$ muestra que la multiplicación por $q^u$ es invertible; por tanto no hay elementos no nulos anulados por ella. Así,
\[
T_q\cong\bigoplus_{j=1}^{v}R/(q^{e_j}).
\]
El cociente $T_q/qT_q$ es un espacio vectorial sobre el campo $k_q=R/(q)$, y cada factor anterior aporta una copia de $k_q$. En particular, $\dim_{k_q}(T_q/qT_q)=v$ recupera el número de factores.

Para recuperar los exponentes, usamos todas las \emph{capas sucesivas}
\[
d_\ell=\dim_{k_q}\left(q^{\ell-1}T_q/q^\ell T_q\right),
\qquad \ell\geq1.
\]
Estas capas sí están anuladas por $q$. En un factor $R/(q^e)$, la capa es cero si $e<\ell$; si $e\geq\ell$, el homomorfismo
\[
R\longrightarrow q^{\ell-1}(R/(q^e))/q^\ell(R/(q^e)),
\qquad a\longmapsto aq^{\ell-1}+q^\ell(R/(q^e))
\]
es sobreyectivo y tiene núcleo $(q)$: la pertenencia del representante a $(q^\ell)$, módulo $(q^e)$, equivale a $q\mid a$, cancelando $q^{\ell-1}$ en el dominio. Por ello la capa es una copia de $k_q$. Concluimos
\[
d_\ell=\#\set{j\given e_j\geq\ell},\qquad
d_\ell-d_{\ell+1}=\#\set{j\given e_j=\ell}.
\]
Los números intrínsecos $d_\ell$ determinan el multiconjunto completo de exponentes. Esto demuestra la unicidad de los divisores elementales.

Finalmente recuperamos los factores invariantes. Para cada primo, ordenamos sus exponentes de menor a mayor. El número $m$ de factores invariantes no unidades es la mayor cantidad de factores elementales que aparece para un primo: no puede superarla, y un primo que divide a $a_1$ divide a todos los $a_j$, por lo que la alcanza. Completamos cada lista de exponentes hasta longitud $m$ añadiendo ceros al principio, y multiplicamos las potencias alineadas en cada posición. La cadena $a_1\mid\cdots\mid a_m$ fuerza exactamente esa alineación, así que determina cada ideal $(a_j)$. Si $T=0$, $m=0$ y no hay listas que reconstruir. Queda probada toda la unicidad.
\end{proof}

\begin{nota}[Corrección del cociente en la página 10]
En la parte inferior se utiliza $T_q/q^\ell T_q$ como si fuera un espacio vectorial sobre $R/(q)$ para todo $\ell$. En general no lo es: en $\mathbb Z/4\mathbb Z$ la multiplicación por $2$ no es cero, aunque el módulo está anulado por $2^2$. La corrección consiste en tomar $q^{\ell-1}T_q/q^\ell T_q$; su dimensión cuenta exponentes al menos $\ell$, como se acaba de demostrar. Se conserva la observación manuscrita válida
\[
q^i\bigl(R/(q^e)\bigr)\cong R/(q^{e-i})\quad(0\leq i<e),
\]
mediante $a\mapsto aq^i$, cuyo núcleo es $(q^{e-i})$; para $i\geq e$ la imagen es cero. La reconstrucción final de los factores invariantes explicita el paso de unicidad que no se desarrolla en las pizarras.
\end{nota}

\newpage
%%%%%%%%%%%%%%%%%%%%%%%%%%%%%%%%%%%%%%%%%%%%%%%%%%%%%%%%%%%%%%%%%%%%%%%%%%%%%%%
\section{Un ejemplo de presentación y forma de Smith}
%%%%%%%%%%%%%%%%%%%%%%%%%%%%%%%%%%%%%%%%%%%%%%%%%%%%%%%%%%%%%%%%%%%%%%%%%%%%%%%
% P11.
Sea ahora $R=k[x]$, con $k$ un campo, y sea el módulo presentado por
\[
A=\begin{pmatrix}x&0\\0&x-1\end{pmatrix}:R^2\to R^2,
\qquad M=\coker A=R^2/\im A.
\]
La diagonal inicial no es una cadena de divisibilidad, pues $x\nmid x-1$. Las operaciones elementales por filas y columnas corresponden a multiplicar por matrices invertibles y conservan el conúcleo salvo isomorfismo. Denotando las filas por $L_i$ y las columnas por $C_i$, reproducimos todos los pasos de la pizarra:
\begin{align*}
\begin{pmatrix}x&0\\0&x-1\end{pmatrix}
&\xrightarrow{C_1\leftarrow C_1+C_2}
\begin{pmatrix}x&0\\x-1&x-1\end{pmatrix}
\xrightarrow{L_2\leftarrow L_2-L_1}
\begin{pmatrix}x&0\\-1&x-1\end{pmatrix}\\[3pt]
&\xrightarrow{L_2\leftarrow-L_2}
\begin{pmatrix}x&0\\1&1-x\end{pmatrix}
\xrightarrow{L_1\leftrightarrow L_2}
\begin{pmatrix}1&1-x\\x&0\end{pmatrix}\\[3pt]
&\xrightarrow{C_2\leftarrow C_2-(1-x)C_1}
\begin{pmatrix}1&0\\x&-x(1-x)\end{pmatrix}
\xrightarrow{L_2\leftarrow L_2-xL_1}
\begin{pmatrix}1&0\\0&x^2-x\end{pmatrix}.
\end{align*}
La última diagonal cumple $1\mid x^2-x$ y es una \emph{forma normal de Smith} de $A$. En consecuencia,
\[
M\cong R/(1)\oplus R/(x^2-x)\cong R/(x^2-x),
\]
que exhibe su único factor invariante no trivial. Como $x^2-x=x(x-1)$ y los ideales $(x)$ y $(x-1)$ son comaximales, la descomposición en divisores elementales es
\[
M\cong R/(x)\oplus R/(x-1).
\]
El argumento funciona en cualquier característica, pues $x-(x-1)=1$. Para justificar la conservación del módulo presentado, si $A'=PAQ$ con $P,Q$ invertibles, entonces $\im A'=P(\im A)$; el automorfismo $P$ del codominio induce el isomorfismo de cocientes. Así se completa la interpretación de las operaciones utilizadas.

\begin{obs}[Orientación bibliográfica]
Las páginas 1--7 de \citet{aluffi2009} reproducen VI.5.2 (recordatorio), VI.5.1, VI.5.3--6 y sus demostraciones manuscritas: submódulos, bases compatibles, resoluciones, rango y existencia de la clasificación. Se resuelve la remisión al ejercicio VI.5.6. Las páginas 8--10 contienen las definiciones VI.4.4/6, la comprobación remitida a VI.4.4, el lema VI.5.7 y el desarrollo manuscrito de unicidad por anuladores y componentes primarias. La página 11 desarrolla el ejemplo de Smith sin localizador numerado visible. La alusión a «ejercicio 1, tarea 4» se concreta aquí en las dos afirmaciones efectivamente usadas: producto igual a intersección de ideales comaximales y teorema chino del resto; no se atribuye a esa remisión un enunciado adicional. Las pruebas añadidas y la corrección de las capas se indican en el texto. La consulta se limita a los fragmentos originales.
\end{obs}
